\section{Simultaneous driven hotspots in an invariant Maxwell geometry}
\label{hs:section}

\noindent\textbf{Status and scope.}
Theorem~\ref{hs:packing} proves a driven, two-dimensional cross-sectional
version of hotspot packing, with physical incoming power per unit invariant
length.  It uses ordinary binary dielectrics at one fixed positive frequency,
one incident field, and smooth rounded interfaces.  The proof includes local
finite-frequency excitation, global illumination, and the rounding limit.
It does not prove packing of finite-three-dimensional point hotspots or
packing by a normalized resonator mode.  Those distinctions are essential
when comparing with the modal question in \cite{millerMonticone2026}.

\subsection{Physical reduction, normalization, and the exact target}
\label{hs:setup}

Fix $k>0$, $\omega=ck$, and $\varepsilon_+>1$.  Set
$b=\varepsilon_+^{-1}\in(0,1)$ and $Z_0=\sqrt{\mu_0/\varepsilon_0}$.
Every coefficient below is independent of $z$.  Put $a=1/\varepsilon_r$,
so that $b\leq a\leq1$, and consider
\begin{equation}
 H=(0,0,u),\qquad
 E=\frac{ia}{\omega\varepsilon_0}(\partial_yu,-\partial_xu,0),
 \qquad L_a u:=\nabla\cdot(a\nabla u)+k^2u=0.
 \label{hs:tm}
\end{equation}
These are the full three-dimensional Maxwell equations restricted to
$z$-invariant fields, rather than an electrostatic approximation.
Indeed $\nabla\times H=-i\omega\varepsilon_0\varepsilon_rE$, and
$\nabla\times E=i\omega\mu_0H$ is exactly $L_au=0$.
The weak scalar transmission conditions are $[u]=[a\partial_nu]=0$.
They give continuity of tangential $H$ and $E$, and of normal displacement.
The local electric energy is finite precisely when $\nabla u$ is locally
square integrable, since
\begin{equation}
 \int_Q\varepsilon_0\varepsilon_r|E|^2\,dx\,dy
 =\frac{1}{\omega^2\varepsilon_0}\int_Q a|\nabla u|^2\,dx\,dy.
 \label{hs:energy}
\end{equation}

The incident field is a smooth Herglotz wave
\begin{equation}
 u_g^{\rm in}(x)=\left(\frac{k}{\pi Z_0}\right)^{1/2}
   \int_{S^1}g(d)e^{ikd\cdot x}\,ds(d),\qquad g\in C^\infty(S^1).
 \label{hs:herglotz}
\end{equation}
The scattered field satisfies the two-dimensional outgoing Sommerfeld
condition.  All material lies in a bounded cross section, and the exterior
is vacuum.  This convention specifies the incoming physical channels;
no freely chosen homogeneous outgoing solution is added.
Stationary phase, for $x=r\widehat x$, gives
\begin{equation}
 u_g^{\rm in}(r\widehat x)=
 \left(\frac{2}{Z_0r}\right)^{1/2}
 \left[e^{i(kr-\pi/4)}g(\widehat x)
       +e^{-i(kr-\pi/4)}g(-\widehat x)\right]+O(r^{-3/2}).
 \label{hs:asymptotic-input}
\end{equation}
An incoming radial wave has inward flux $Z_0|u|^2/2$; an outgoing
one has outward flux $Z_0|u|^2/2$.  Therefore
\begin{equation}
 P_{\rm in}\text{ per unit }z\text{ length}=\|g\|_{L^2(S^1)}^2.
 \label{hs:power}
\end{equation}
This is incoming, rather than net, power.  The latter vanishes for a lossless
scatterer after the outgoing power is included.  A plane wave alone would
not have this finite incoming-power normalization.

\begin{definition}[Localized peaks used here]
\label{hs:peak-definition}
For a smooth-interface, piecewise homogeneous configuration, let
$E_*(x)$ equal $|E(x)|$ off interfaces and the maximum of the finitely
many one-sided values at an interface.  A collection of $N$
localized peaks at length target $\ell>0$ and amplitude target $M>0$
means $N$ pairwise disjoint closed disks $Q_j$, all of diameter below
$\ell$, such that
\begin{equation}
 P_j:=\max_{\overline Q_j,\,\text{one-sided}}|E|>M,
 \qquad
 \max_{\partial Q_j,\,\text{one-sided}}|E|<P_j/2.
 \label{hs:peak-test}
\end{equation}
We will in fact bound the field on an annular collar of each boundary.
Consequently a connected component of the superlevel set
$\{E_*>P_j/2\}$ containing a maximizing point stays in $Q_j$.
Its diameter is below $\ell$.  The corresponding intensity half-maximum
component is contained in it, since $P_j/\sqrt2>P_j/2$.
No positive lower bound on the energy in each peak is stipulated.
\end{definition}

\begin{theorem}[PROVED: simultaneous rounded driven hotspot packing]
\label{hs:packing}
For every fixed $k>0$ and $\varepsilon_+>1$, every nonempty open
patch $\mathcal P\subset\mathbb R^2$, every finite $N\geq1$, and every
$\ell,M>0$, there exist $2N$ bounded dielectric inclusions with pairwise
disjoint closures and $C^\infty$ boundaries, all contained in
$\mathcal P$, and one $g\in C^\infty(S^1)$ with $\|g\|_2=1$, such that:
\begin{enumerate}
 \item $\varepsilon_r=\varepsilon_+$ inside the inclusions and
       $\varepsilon_r=1$ outside them;
 \item the unique outgoing scattering solution of \eqref{hs:tm}
       driven by \eqref{hs:herglotz} has $N$ localized peaks in
       $\mathcal P$ satisfying Definition~\ref{hs:peak-definition};
 \item the same one incident field excites every peak, at the same
       frequency and with unit incoming power per unit $z$ length.
\end{enumerate}
There is no uniform lower bound on the rounding radius, on any peak's
integrated energy, or on any singular coefficient as $N$ or the target
parameters change.  The physical construction is infinitely extruded
in $z$ and is not a finite-volume three-dimensional device.
\end{theorem}

The strict contrast assumption matters: for $\varepsilon_+=1$ the material
is vacuum.  Cauchy--Schwarz in \eqref{hs:herglotz} then gives a finite
pointwise electric-field bound at fixed unit incoming power and fixed $k$,
so the arbitrary-amplitude conclusion cannot include that endpoint.

\subsection{A local singular projection at positive frequency}
\label{hs:local}

Corner and edge singularities of positive-coefficient Maxwell transmission
problems are established in \cite{costabel1999}.  Here we give the needed
two-dimensional construction directly.  In particular, we do not transfer
a bounded-domain boundary condition from that paper into an unproved
scattering claim.

\begin{lemma}[PROVED: a binary checkerboard angular exponent]
\label{hs:angular}
Let $a(\theta)$ alternate between $1$ and $b$ in successive sectors of
angle $\pi/2$.  There is a nonzero real periodic $\phi\in H^1(S^1)$ such
that
\begin{equation}
 -(a\phi')'=\alpha^2a\phi,
 \qquad \alpha=\frac4\pi\arctan\sqrt b\in(0,1),
 \qquad \int_0^{2\pi}a\phi^2\,d\theta=1.
 \label{hs:angular-eigenproblem}
\end{equation}
Moreover $\int a\phi=0$ and $\int a|\phi'|^2=\alpha^2$.
\end{lemma}

\begin{proof}
In a sector with constant coefficient $a_0$, the state
$(\phi,a_0\phi'/\alpha)^T$ propagates through angle $L$ by
\begin{equation}
 M_{a_0}(L)=
 \begin{pmatrix}
 \cos\alpha L & a_0^{-1}\sin\alpha L\\
 -a_0\sin\alpha L & \cos\alpha L
 \end{pmatrix}.
 \label{hs:transfer}
\end{equation}
Both state components are continuous at an interface.  The half-turn
matrix $B=M_b(\pi/2)M_1(\pi/2)$ has determinant one and trace
$2\cos^2t-(b+b^{-1})\sin^2t$, where $t=\alpha\pi/2$.
With $t=2\arctan\sqrt b$ its trace is $-2$.  Choose a nonzero
eigenvector of $B$ with eigenvalue $-1$.  The resulting angular function
is antiperiodic after a half turn and periodic after a full turn.
It satisfies the stated weak eigenproblem.  Integrating that equation
and then testing it against $\phi$ proves the two identities.
\end{proof}

\begin{lemma}[PROVED: continuous finite-frequency singular coefficient]
\label{hs:coefficient}
Suppose the checkerboard coefficient of Lemma~\ref{hs:angular} holds in
$B_R(0)$ and $u\in H^1(B_R)$ weakly satisfies $L_au=0$ there.
For almost every $r$, put
\begin{equation}
 U(r)=\int_0^{2\pi}a(\theta)u(r,\theta)\phi(\theta)\,d\theta,
 \qquad F(r)=\int_0^{2\pi}u(r,\theta)\phi(\theta)\,d\theta.
 \label{hs:projections}
\end{equation}
There is a complex number $c(u)$, linear in $u$, such that
\begin{align}
 U(r)&=c(u)r^\alpha+p(r),\notag\\
 p(r)&=-\frac{k^2}{2\alpha}\left[
 r^\alpha\int_0^r s^{1-\alpha}F(s)\,ds
 -r^{-\alpha}\int_0^r s^{1+\alpha}F(s)\,ds\right].
 \label{hs:exact-projection}
\end{align}
In particular,
\begin{equation}
 |p(r)|\leq C_{b,\phi,\alpha}k^2r\|u\|_{L^2(B_r)},
 \qquad
 |c(u)|\leq C_{R,k,b,\phi}\|u\|_{L^2(B_R)}.
 \label{hs:projection-estimates}
\end{equation}
Thus $c$ is continuous even for the local $L^2$ topology on the
space of solutions.  If $c(u)\ne0$, ordinary points strictly inside the
homogeneous sectors have arbitrarily large $|E|$ as they approach the
origin.
\end{lemma}

\begin{proof}
Testing the weak equation with a product of a radial test function and
$\phi$ gives the distributional ordinary differential equation
\begin{equation}
 (rU')'-\frac{\alpha^2}{r}U=-k^2rF.
 \label{hs:radial-ode}
\end{equation}
Variation of constants gives the particular solution in
\eqref{hs:exact-projection}.  The two integrals converge because
$F\in L^2((0,R),r\,dr)$ and $\alpha<1$.  For example,
\begin{equation}
 \int_0^r s^{1-\alpha}|F(s)|\,ds
 \leq\left(\int_0^r s|F(s)|^2\,ds\right)^{1/2}
       \frac{r^{1-\alpha}}{\sqrt{2-2\alpha}}.
 \label{hs:weighted-integral}
\end{equation}
The second integral is bounded in the same way, with exponent
$1+\alpha$ and denominator $\sqrt{2+2\alpha}$.  These estimates prove
the first bound in \eqref{hs:projection-estimates}.  With
$I_\mp(r)=\int_0^r s^{1\mp\alpha}F(s)\,ds$, differentiation gives
\begin{equation}
 p'(r)=-\frac{k^2}{2}
       \left[r^{\alpha-1}I_-(r)+r^{-\alpha-1}I_+(r)\right],
 \qquad |p'(r)|\leq Ck^2\|u\|_{L^2(B_r)}.
 \label{hs:particular-derivative}
\end{equation}
Thus $p'$ is locally square integrable with radial measure.
The general remaining homogeneous term is $c r^\alpha+d r^{-\alpha}$.
Since $u\in H^1$, $U'$ is square integrable with measure $r\,dr$,
which excludes $d\ne0$.

To avoid any assumption about point evaluation of an $H^1$ function,
solve \eqref{hs:exact-projection} for $c$ and average over
$r\in(R/2,R)$ with the probability measure $2\,dr/R$.
The resulting formula is
\begin{equation}
 c(u)=\frac2R\int_{R/2}^R
 \left\{r^{-\alpha}U(r)+\frac{k^2}{2\alpha}
 \left[\int_0^r s^{1-\alpha}F(s)\,ds
 -r^{-2\alpha}\int_0^r s^{1+\alpha}F(s)\,ds\right]\right\}\,dr.
 \label{hs:coefficient-functional}
\end{equation}
Cauchy--Schwarz and \eqref{hs:weighted-integral} bound this expression
by $C\|u\|_{L^2(B_R)}$.  This proves continuity without invoking a
weighted corner-expansion theorem.

For almost every radius, angular integration by parts and
\eqref{hs:angular-eigenproblem} yield
\begin{equation}
 \alpha|U(r)|\leq
 \left(\int_0^{2\pi}a|\partial_\theta u(r,\theta)|^2\,d\theta\right)^{1/2}.
 \label{hs:gradient-lower-bound}
\end{equation}
Since $p(r)=o(r)$ and $\alpha<1$, $c\ne0$ implies
$|U(r)|\geq |c|r^\alpha/2$ for sufficiently small radii.
Consequently the essential supremum of $|\nabla u|$ on each such
circle is at least $\alpha|c|r^{\alpha-1}/(2\sqrt{2\pi})$.
Here $a\leq1$ was used.  Interior analyticity in each constant-coefficient
sector permits the large values to be selected at ordinary points,
off the interfaces.  Formula~\eqref{hs:tm} gives
$|E|\geq b|\nabla u|/(\omega\varepsilon_0)$ there.
\end{proof}

\begin{remark}[Energy and what has actually been expanded]
\label{hs:integrability}
The static singular scalar field $u_0=r^\alpha\phi$ has
\begin{equation}
 \int_{B_R}a|\nabla u_0|^2\,dx\,dy=\alpha R^{2\alpha}.
 \label{hs:static-energy}
\end{equation}
Its electric-field exponent is $\alpha-1$, and its energy is integrable.
The finite-frequency fields above have finite local energy by membership
in $H^1$.  Formula~\eqref{hs:exact-projection} is an exact singular
\emph{angular projection}; no unproved pointwise expansion of all vector
components is needed for the lower-bound argument.  In particular,
finite local energy does not bound the peak electric field.
\end{remark}

\begin{lemma}[PROVED: a local finite-frequency solution with $c\ne0$]
\label{hs:local-excitation}
For every fixed $k>0$ there is $R_*>0$, depending on $k$ and $b$, such
that for $0<R<R_*$ the checkerboard equation in $B_R$ has an $H^1$
solution with nonzero coefficient in Lemma~\ref{hs:coefficient}.
\end{lemma}

\begin{proof}
Scale $x=Ry$ and set $\kappa=kR$.  In the unit disk solve
$\nabla\cdot(a\nabla v)+\kappa^2v=0$ with Dirichlet trace $\phi$.
The static extension $v_0=|y|^\alpha\phi$ belongs to $H^1$ and has
that trace.  Write $v=v_0+w$, $w\in H_0^1(B_1)$.  Its equation is
\begin{equation}
 \int_{B_1}a\nabla w\cdot\nabla\overline\psi
       -\kappa^2w\overline\psi
 =\kappa^2\int_{B_1}v_0\overline\psi,
 \qquad \psi\in H_0^1(B_1).
 \label{hs:dirichlet-perturbation}
\end{equation}
If $C_P$ is the unit-disk Poincare constant, choose
$\kappa^2C_P^2<b/2$.  The left form is coercive; Lax--Milgram gives
$\|w\|_{H^1}\leq C_b\kappa^2\|v_0\|_2$ and hence a uniform bound
on $\|v\|_2$ for these small $\kappa$.
Because the boundary trace has $U(1)=\int a\phi^2=1$,
\eqref{hs:exact-projection} at $r=1$ gives
\begin{equation}
 c(v)=1+\frac{\kappa^2}{2\alpha}
 \left[\int_0^1s^{1-\alpha}F_v(s)\,ds
       -\int_0^1s^{1+\alpha}F_v(s)\,ds\right]
       =1+O_b(\kappa^2).
 \label{hs:nonzero-coefficient}
\end{equation}
The integral estimates already proved justify the last bound.
Reduce $R_*$ until $|c(v)-1|<1/2$.
For $u(x)=v(x/R)$ the physical coefficient is $R^{-\alpha}c(v)$,
which is nonzero.  The frequency remained the prescribed positive $k$;
only the local matching disk was made small.
\end{proof}

\subsection{Outgoing well-posedness and global excitation}
\label{hs:global}

We spell out the scalar scattering and Runge arguments because density
of free-space Herglotz waves alone would not establish excitation of
the total field in a material.  General Herglotz approximation results
such as \cite{ochs1989} motivate the method, but are not substituted for
the following transmission argument.

\begin{lemma}[PROVED: scattering for the finite polygonal geometry]
\label{hs:wellposed}
Let $a$ be real, take values in $[b,1]$, equal one outside a disk, and
be constant on a finite polygonal partition.  The outgoing scattering
problem for \eqref{hs:tm} has a unique $H^1_{\rm loc}$ total field for
each incident Herglotz wave.  The map from $g\in L^2(S^1)$ to $u_g$ in
$H^1$ on any fixed bounded disk is continuous.  An outgoing solution
also exists uniquely for any compactly supported $L^2$ forcing.
\end{lemma}

\begin{proof}
Choose $B_L$ enclosing the material with a vacuum collar.  Let
$\Lambda_L:H^{1/2}(\partial B_L)\to H^{-1/2}(\partial B_L)$ be the
exact exterior outgoing Dirichlet-to-Neumann map.  Its Fourier
multiplier on angular mode $m\in\mathbb Z$ is
\begin{equation}
 \lambda_m=k\frac{(H_{|m|}^{(1)})'(kL)}{H_{|m|}^{(1)}(kL)}.
 \label{hs:dtn}
\end{equation}
The Hankel denominator is nonzero on the positive real axis, since
the real Bessel functions $J_m,Y_m$ have nonzero Wronskian.
Their large-order asymptotics give
$\lambda_m=-|m|/L+O((1+|m|)^{-1})$ as $|m|\to\infty$.
Thus the form
\begin{equation}
 B_a(v,w)=\int_{B_L}\left(a\nabla v\cdot\nabla\overline w
                           -k^2v\overline w\right)
        -\langle\Lambda_Lv,w\rangle_{\partial B_L}
 \label{hs:form}
\end{equation}
satisfies, by the trace inequality,
\begin{equation}
 \operatorname{Re}B_a(v,v)+C\|v\|_{L^2(B_L)}^2
       \geq c\|v\|_{H^1(B_L)}^2.
 \label{hs:garding}
\end{equation}
The constants can be chosen using only $b,k,L$, not the positions or
curvatures of the interfaces: a finite upper bound for the real parts
of the multipliers suffices, and
$\|v\|_{L^2(\partial B_L)}^2\leq
\epsilon\|v\|_{H^1(B_L)}^2+C_\epsilon\|v\|_2^2$ absorbs it.
The compact embedding $H^1(B_L)\hookrightarrow L^2(B_L)$ makes this
a Fredholm form of index zero.

For the homogeneous outgoing problem, take the imaginary part of
Green's identity.  The far-field outgoing power is nonnegative and
must vanish.  Equivalently, an exterior Hankel expansion and its
Wronskians show that all outgoing coefficients vanish.  The solution
is zero in the exterior vacuum.  It continues as zero across every
open polygon edge: its value and its conormal derivative are zero
on that edge, so extension by zero to the other side solves that
side's constant-coefficient Helmholtz equation in a small ball.
Interior analytic uniqueness forces it to vanish there too.
A path can avoid the finitely many vertices, and cross only open
edges.  Propagation along such paths proves that the solution is zero
everywhere.  This argument also proves unique continuation through
this particular partition in any connected open subset.
Fredholm alternative gives existence and a bounded inverse.

For scattering, the right-hand side of \eqref{hs:form} is the boundary
functional with density
$\partial_nu_g^{\rm in}-\Lambda_Lu_g^{\rm in}$.
For a volume source $L_aw=f$, it is $-\int f\overline v$.
Both are continuous functionals.  The incident field and all its
derivatives on bounded sets are bounded by $C\|g\|_2$, directly from
\eqref{hs:herglotz}.  The bounded inverse gives the asserted continuity.
\end{proof}

\begin{lemma}[PROVED: local Runge approximation by total scattering fields]
\label{hs:runge}
Under Lemma~\ref{hs:wellposed}, let $Q$ be a closed disk whose
complement is connected.  A solution of $L_av=0$ in an open
neighborhood of $Q$ can be approximated in $L^2(Q)$ by restrictions
of total fields $u_g$, with $g\in C^\infty(S^1)$.
The disk may cross material interfaces and contain a checkerboard
vertex.  The same statement holds for a finite disjoint union of
disks if its complement is connected.
\end{lemma}

\begin{proof}
We use complex bilinear pairings, with conjugation absorbed into the
annihilating function.  Suppose $f\in L^2(Q)$ satisfies
$\int_Q f u_g=0$ for all smooth $g$.  Extend $f$ by zero and solve
$L_aw=f$ outgoing.  For a plane-wave incidence $e^{ikd\cdot x}$,
write $u_d=e^{ikd\cdot x}+u_d^{\rm sc}$ for its total field.
Linearity of scattering and Fubini imply
$\int_Qf u_d=0$ for all $d\in S^1$: the quantity is continuous in
$d$ and its pairing with every smooth density vanishes.

Green's second identity on a sufficiently large circle gives
\begin{equation}
 \int_Q f u_d
 =\int_{\partial B_R}
       (\partial_nw\,u_d-w\,\partial_nu_d)\,ds.
 \label{hs:reciprocity-boundary}
\end{equation}
The part pairing $w$ with $u_d^{\rm sc}$ is zero.  For example, expand
both in outgoing Hankel modes; bilinear angular orthogonality pairs
$m$ with $-m$, whose equal radial Hankel functions have zero
Wronskian.  The remaining part is the free-space incident wave paired
with $w$.  The exterior Green representation with the fundamental
solution $(i/4)H_0^{(1)}(k|x-y|)$ identifies this as a nonzero constant
times $w_\infty(-d)$.  This can also be checked by the
Jacobi--Anger expansion of $e^{ikd\cdot x}$ and the $J_m,H_m^{(1)}$
Wronskian.  Hence $w_\infty=0$.  The outgoing Hankel expansion gives
$w=0$ outside a sufficiently large circle.

Now $L_aw=0$ in $\mathbb R^2\setminus Q$.  That set is connected,
and the edgewise continuation proved in Lemma~\ref{hs:wellposed}
implies $w=0$ throughout it.  Thus $w\in H_0^1(Q)$ and both its
Dirichlet and conormal traces on $\partial Q$ vanish.  The latter
statement is in the weak $H^{-1/2}$ sense: the global equation has
only the volume source $f$, and no surface distribution.  Applying
weak Green's identity to any local solution $v$ gives
\begin{equation}
 \int_Qfv
 =\int_Q(L_aw)v
 =\int_Qw(L_av)=0.
 \label{hs:annihilator}
\end{equation}
For a union of disks apply the same identity on each component.
The Hahn--Banach annihilator characterization of a closed subspace
of $L^2(Q)$ now proves the density assertion.  No interior
transmission-eigenvalue exception was introduced: the argument used
the invertible outgoing source problem, not an inverse of an
interior transmission operator.
\end{proof}

\begin{proposition}[PROVED: one incident field excites every selected corner]
\label{hs:simultaneous}
In a fixed finite polygonal geometry with $N$ disjoint checkerboard
disks of the same contrast, each local coefficient $c_j(u_g)$ is a
nonzero continuous linear functional of $g$.  There exists one
$g\in C^\infty(S^1)$ with $\|g\|_2=1$ for which
$c_j(u_g)\ne0$ for every $j$.
\end{proposition}

\begin{proof}
Continuity follows from Lemmas~\ref{hs:coefficient} and
\ref{hs:wellposed}.  Inside corner $j$ choose a sufficiently small
disk, and then a concentric disk twice its radius still within the
checkerboard region.  Lemma~\ref{hs:local-excitation} on the larger
disk supplies a local finite-frequency solution with $c_j\ne0$.
Lemma~\ref{hs:runge} approximates it on the smaller disk by total
scattering fields.  The $L^2$ continuity of $c_j$ shows that at least
one smooth density $g_j$ has $c_j(u_{g_j})\ne0$.

Let $g(t)=\sum_{m=1}^N t^{m-1}g_m$.  For each fixed $j$,
$c_j(u_{g(t)})$ is a polynomial in $t$ with nonzero coefficient of
$t^{j-1}$.  It has only finitely many zeros.  Choose a real $t$
outside the union of these finite root sets.  Then $g(t)\ne0$ and
all the coefficients are nonzero.  Dividing by $\|g(t)\|_2$ preserves
that property and enforces \eqref{hs:power}.  This finite-union
argument requires no uniform quantitative lower bound on the
coefficients and uses exactly one illumination.
\end{proof}

\subsection{Rounding without an assumed rounding power law}
\label{hs:rounding}

\begin{lemma}[PROVED: outgoing solutions converge under local rounding]
\label{hs:rounding-convergence}
Let $a$ be as in Lemma~\ref{hs:wellposed}.  Let $a_\rho$ equal one
in a common exterior collar, satisfy $b\leq a_\rho\leq1$, and
converge to $a$ almost everywhere.  Assume each $a_\rho$ defines a
well-posed outgoing problem, as do the finite disjoint smooth binary
inclusions used below.  For a fixed incident field,
\begin{equation}
 u_\rho\longrightarrow u\quad\hbox{strongly in }H^1(B_L).
 \label{hs:h1-convergence}
\end{equation}
If $a_\rho=a$ in a fixed neighborhood of a compact observation set
away from polygon vertices, convergence there is $C^1$ on each
closed material side.  The compact set may cross a straight
interface.  Thus both ordinary field values and one-sided field
bounds on fixed observation circles converge.
\end{lemma}

\begin{proof}
All forms $B_{a_\rho}$ satisfy the same G\aa rding estimate
\eqref{hs:garding}.  Their inverse norms on
$H^1(B_L)^*\to H^1(B_L)$ are uniformly bounded for sufficiently
small $\rho$.  Otherwise there would exist $v_\rho$ with
$\|v_\rho\|_{H^1}=1$ and
$\|B_{a_\rho}v_\rho\|_{H^1{}^*}\to0$.
Pass to a weakly convergent $H^1$ subsequence, strongly convergent
in $L^2$.  For every fixed test $z$, dominated convergence gives
$(a_\rho-a)\nabla z\to0$ in $L^2$, so the limit $v$ solves
$B_av=0$.  Lemma~\ref{hs:wellposed} implies $v=0$.
Estimate~\eqref{hs:garding} then forces
$\|v_\rho\|_{H^1}\to0$, a contradiction.

Subtracting the equations with the same boundary forcing gives
\begin{equation}
 B_{a_\rho}(u_\rho-u,z)
 =\int_{B_L}(a-a_\rho)\nabla u\cdot\nabla\overline z.
 \label{hs:rounding-difference}
\end{equation}
The right-hand functional tends to zero in $H^1{}^*$ because
$(a-a_\rho)\nabla u\to0$ in $L^2$ by dominated convergence.
The uniform inverse bound proves \eqref{hs:h1-convergence}.

For completeness, the additional local regularity needed here does
not rely on estimates uniform near the rounded tips.  In a ball
with constant coefficient, interior elliptic estimates applied to
$u_\rho-u$ give convergence of all derivatives on smaller balls.
In a ball meeting one unchanged straight interface, tangential
difference quotients preserve both transmission conditions.
Caccioppoli estimates for these quotients, iterated to each order,
bound tangential derivatives by the $H^1$ norm on a larger ball.
On each side the equation
$\partial_n^2v=-\partial_t^2v-(k^2/a_\pm)v$
then bounds the second normal derivative.  Differentiating this
identity on each material side inductively bounds derivatives with
two or more normal factors by the tangential derivatives already
controlled and lower-order derivatives.  This yields, for
each integer $m$ and fixed nested interface balls, a piecewise
$H^m$ estimate with a constant depending on those balls, $b,k,m$.
Sobolev embedding with $m\geq3$ in two dimensions gives the claimed
one-sided $C^1$ convergence.  A finite cover handles the compact
observation set.  None of these balls contains a moving rounded tip.

The well-posedness assertion for the smooth inclusions follows by
the same Fredholm and radiation argument.  For uniqueness, zero
exterior Cauchy data extend by zero across a smooth material boundary
as a weak solution of that inclusion's constant-coefficient
Helmholtz equation, so interior analyticity again applies.  Standard
flattening and the same transmission regularity argument show that
each fixed rounded solution has bounded, continuous one-sided
gradients up to its smooth interfaces.  These last bounds may depend
arbitrarily badly on $\rho$; the proof never assumes otherwise.
\end{proof}

\begin{proof}[Proof of Theorem~\ref{hs:packing}]
Choose distinct points $p_1,\ldots,p_N$ in $\mathcal P$ and small
pairwise disjoint closed squares $p_j+[-s_j,s_j]^2$ contained there.
For the sharp reference configuration, put dielectric material in
the two opposite open squares
\begin{equation}
 p_j+(0,s_j)^2,
 \qquad p_j+(-s_j,0)^2,
 \label{hs:square-geometry}
\end{equation}
and vacuum elsewhere.  The coefficient near $p_j$ is the required
four-sector checkerboard.  The two material squares initially
touch only at $p_j$.  There is no source at this touching point;
the problem is the global uniformly elliptic weak equation.

Choose concentric observation disks $Q_j=B_{r_j}(p_j)$ with
$2r_j<\ell$, their closures inside these reference blocks, and
small fixed annular collars that avoid every polygon vertex.
Proposition~\ref{hs:simultaneous} supplies one smooth unit-norm $g$
for which every $c_j(u_g)$ is nonzero.  Keep this $g$ fixed from
now on.  The sharp solution is piecewise smooth on the collars;
let $B<\infty$ bound all one-sided values of $|E|$ on them.
Choose an electric-field threshold $S>\max\{M,2B\}$.
By Lemma~\ref{hs:coefficient}, for every $j$ one can choose an
ordinary point $x_j\in Q_j$, strictly off the interfaces, with
\begin{equation}
 |E(x_j)|>4S.
 \label{hs:selected-values}
\end{equation}
All these choices precede rounding.

Round every vertex of every material square \emph{inwards} within
a distance $\rho$, leaving all edge portions outside the vertex
balls unchanged.  A smooth curve joining the shortened straight
edges with all endpoint derivatives matched gives $C^\infty$
boundaries.  For $\rho>0$ the two opposite inclusions have disjoint
closures.  The resulting binary coefficient $a_\rho$ differs from
$a$ only in the union of finitely many shrinking vertex balls, so
$a_\rho\to a$ almost everywhere and retains exactly the original
material bounds and exterior background.

Choose one sufficiently small $\rho$ so that no changed vertex
ball meets any selected point or collar.  Lemma~\ref{hs:rounding-convergence}
and finiteness of $N$ then permit a further decrease of this same
$\rho$ until simultaneously
\begin{equation}
 |E_\rho(x_j)|>3S\quad(1\leq j\leq N),
 \qquad
 \sup_{\text{all collars, one-sided}}|E_\rho|<S.
 \label{hs:rounded-values}
\end{equation}
For each fixed rounded geometry the one-sided field is bounded and
continuous on the closure of each smooth material subdomain.
Hence its maximum $P_j$ on $\overline Q_j$ is finite and attained
on at least one side, with $P_j>3S>M$.
The collar values are below $S<P_j/2$.  Each maximum is therefore
inside its disk, and its half-amplitude superlevel component cannot
cross the collar.  This is \eqref{hs:peak-test}, with distinct
components in distinct disks.  The inclusions remain entirely in
$\mathcal P$, and the fixed $g$ still has the same unit incoming
power.  This proves the theorem.
\end{proof}

\subsection{Exact scope, costs, and remaining hotspot questions}
\label{hs:limits}

The order of choices in the proof is
\begin{equation}
 (N,\mathcal P,\ell,k,b)\ ;\ \text{sharp geometry}\ ;\ g\ ;\
 (S,x_1,\ldots,x_N)\ ;\ \rho.
 \label{hs:order-of-limits}
\end{equation}
No interchange of the $N\to\infty$ and $\rho\to0$ limits occurs.
Runge approximation can have arbitrarily large illumination norms
before normalization; after normalization the nonzero coefficients
can be arbitrarily small.  This costs a smaller rounding radius,
which is permitted.  The proof supplies no effective rate or
manufacturable feature-size estimate.  It also supplies no lower
bound on the energy fraction in a peak, no independence of different
peaks as communication channels, and no finite spectral bandwidth.
All square dimensions can be chosen inside the given patch, so even
the total \emph{cross-sectional} material support is small.

\begin{proposition}[PROVED: bounded-region energy-to-peak ratios]
\label{hs:regional-ratio}
For the fixed sharp geometry and one illumination used above, along
the rounded sequence the electromagnetic energy on any fixed bounded
cross-sectional disk remains bounded, whereas the electric-field
maximum in every selected corner disk tends to infinity in the
sense that its lower limit is infinite.  Thus the ratio of that
bounded-region energy to the maximum electric energy density in
any one selected disk tends to zero.
\end{proposition}

\begin{proof}
Strong $H^1$ convergence and \eqref{hs:energy} bound the electric
energy; $H=(0,0,u)$ bounds the magnetic energy by $L^2$ convergence.
For each arbitrary threshold, choose the finitely many high ordinary
points in the limiting sharp field first.  Pointwise convergence at
those fixed points then puts all sufficiently small rounding radii
above that threshold.  The weighted density
$\varepsilon_0\varepsilon_r|E|^2$ differs from $|E|^2$ by uniformly
bounded positive factors.  This proves the ratio statement.
\end{proof}

This ratio is deliberately not called a modal volume.  A nonzero
monochromatic scattering field has infinite all-space energy.  A
fixed bounded integration region is a different normalization from
the all-space energy of a mode.  The bounds and examples in
\cite{zhao2020} concern dielectric-resonator mode volumes with their
specified confinement and feature constraints; they cannot replace
the missing modal argument here.

\begin{remark}[Scope: finite extrusion and modal packing]
\label{hs:open3d}
Two specific bridges remain unproved by this section.
First, truncating the $z$-invariant construction to a finite height
changes the exact scalar reduction and introduces end radiation.
A proof of finite-three-dimensional packing would need a vector
Maxwell local singular coefficient, a corresponding physical
illumination density argument, and stability of that coefficient
under finite-height realization and simultaneous rounding.  Our
two-dimensional proof does not supply a uniform estimate that
permits this truncation.  Second, a modal theorem would need an
actual mode or a stated open-resonator normalization and a proof
that the same normalized mode retains all the selected singular
coefficients through the geometry and rounding limits.  The freely
selected Herglotz illumination cannot be substituted for such a mode.
The finite-extrusion route remains unproved by this section. Section
\ref{hs:three-section} supplies a different finite three-dimensional
driven construction with conical point tips and explicit projection
constraints. The modal bridge remains \textbf{OPEN}. These are
additional statements, not assumptions hidden in Theorem~\ref{hs:packing}.
\end{remark}

\begin{remark}[Scope of material extensions]
\label{hs:material-scope}
The construction is an ordinary, reciprocal, nonmagnetic, scalar,
lossless dielectric example.  It therefore also belongs to any
larger material class that contains $1$ and some fixed
$\varepsilon_h>1$ within its admissible bounds, by replacing
$\varepsilon_+$ with $\varepsilon_h$.  The theorem needs neither
near-zero permittivity nor anisotropy.  A positive-definite tensor
class containing these scalar tensors inherits this existence
example, but no anisotropic Runge or corner theorem is claimed.
\end{remark}
